Las conclusiones se ordenan según los objetivos específicos de la Sección~\ref{sec:objetivos}.

\parrafo{O1: modelo integrado frente al caso base.} La matheurística es más barata que el caso base en las 40 semanas BIG, con un ahorro promedio de 9,9\,\% (\$3,9 millones por semana). La comparación es exigente: el caso base corregido externaliza un 70\,\% menos que el anterior. El ahorro viene sobre todo de rutear mejor, no de contratar más: con casi la misma flota se recorre un 16\,\% menos y se externaliza un 28\,\% menos. El costo de esa eficiencia es un aumento de las entregas atrasadas (de 1,3 a 4,1 por semana), que la organización puede limitar subiendo su penalización. El código corregido hace lo que el informe describe y todas sus soluciones pasan la validación independiente.

\parrafo{O2: cuellos de botella.} El recurso que limita la red es el furgón frío. En 26 de las 40 semanas se usan los 18 disponibles, y en esas semanas se externaliza más. Las solicitudes difíciles son casi siempre las mismas para ambos métodos (253 de 287): insumos y equipos grandes con ventanas de retiro de una hora, que necesitan un furgón o dos vehículos a la vez. Duplicar su costo de externalización casi no cambia cuántas se externalizan.

\parrafo{O3: robustez.} La ventaja del método se mantiene en todos los escenarios evaluados, con ahorros entre 7,3 y 10,2\,\%. Un 20\,\% más de demanda sube el costo un 22 a 24\,\%, más que proporcionalmente, porque no hay furgones para crecer; perder un cuarto de la flota lo sube un 4 a 5\,\%. La congestión sube el costo hasta un 10\,\%,, por lo que los costos reportados son un piso. El costo de externalizar, ambiguo en los datos, cambia el nivel de costo pero no la estructura de la solución.

\parrafo{O4: infraestructura.} Las nueve bases se justifican incluso con su costo duplicado. Talagante es la excepción posible: cerrarla no subió el costo en las dos semanas evaluadas, por lo que su cierre merece evaluarse con más semanas. Los hubs casi no se usan, incluso cuando se ofrecen transferencias para todas las solicitudes transferibles: con estos datos, el desvío y la espera del intercambio no compensan. Si hubiera que invertir, conviene más capacidad en furgones antes que más infraestructura.

\parrafo{O5: escalabilidad.} Pendiente: una semana BIG toma 25 minutos en un computador personal, y las 40 semanas SMALL son la primera tarea del Capítulo~\ref{cap:pasos}.
